\subsection{Cache 模拟器设计与实现}

在获得 NEMU 输出的内存访问序列 trace.log 后，本实验进一步编写 Cache 模拟器，对 trace 中记录的访存地址进行逐条模拟访问，从而统计不同 Cache 配置下的 Hit 次数、Miss 次数和命中率。

Cache 模拟器的核心目标并不是模拟完整处理器，而是抽象出 Cache 的基本工作机制，包括地址拆分、组索引计算、Tag 匹配、命中判断以及缺失后的替换过程。通过这种方式，可以在软件层面观察不同 Cache 参数对程序访存性能的影响。

\subsubsection{模拟器输入与输出}

Cache 模拟器的输入为 NEMU 生成的 trace.log 文件。该文件记录了程序运行过程中的访存行为，模拟器从中提取访存地址，并按照给定 Cache 配置依次进行访问模拟。

模拟器的输出主要包括：

\begin{itemize}
    \item 总访存次数；
    \item Cache Hit 次数；
    \item Cache Miss 次数；
    \item Cache 命中率；
    \item 不同配置下的统计结果对比。
\end{itemize}

其中，命中率的计算公式为：

\[
HitRate = \frac{Hit}{Hit + Miss} \times 100\%
\]

\subsubsection{Cache 参数设计}

本实验中的 Cache 模拟器支持对以下参数进行调整：

\begin{itemize}
    \item \texttt{cache\_size}：Cache 总容量；
    \item \texttt{block\_size}：Cache 块大小；
    \item \texttt{associativity}：相联度，即每组包含的 Cache Line 数；
    \item \texttt{replacement\_policy}：替换算法，包括 FIFO 和 LRU；
    \item \texttt{set\_num}：Cache 组数。
\end{itemize}

其中，Cache 组数由 Cache 总容量、块大小和相联度共同决定：

\[
set\_num = \frac{cache\_size}{block\_size \times associativity}
\]

通过改变 \texttt{associativity} 的取值，可以统一表示不同的映射方式：

\begin{itemize}
    \item 当 \texttt{associativity = 1} 时，Cache 为直接映射；
    \item 当 \texttt{associativity = k} 时，Cache 为 $k$ 路组相联；
    \item 当 \texttt{set\_num = 1} 时，Cache 为全相联映射。
\end{itemize}

\subsubsection{地址拆分与映射过程}

对于一条访存地址 \texttt{addr}，模拟器首先根据块大小计算其所属的主存块号：

\[
block\_addr = \left\lfloor \frac{addr}{block\_size} \right\rfloor
\]

随后根据 Cache 组数计算组索引：

\[
set\_index = block\_addr \bmod set\_num
\]

最后计算 Tag：

\[
tag = \left\lfloor \frac{block\_addr}{set\_num} \right\rfloor
\]

因此，每一次访存都可以被映射到某一个 Cache Set 中。模拟器只需要在该组内查找是否存在相同 Tag 的有效 Cache Line，即可判断本次访问是否命中。

\subsubsection{命中检查机制}

模拟器处理每一条访存记录时，首先计算该地址对应的 \texttt{set\_index} 和 \texttt{tag}，然后在对应的 Cache Set 中顺序查找：

\begin{enumerate}
    \item 若存在 \texttt{valid = 1} 且 \texttt{tag} 相同的 Cache Line，则本次访问为 Hit；
    \item 若不存在匹配项，则本次访问为 Miss；
    \item 若发生 Miss 且该组仍有空闲 Cache Line，则直接装入新块；
    \item 若发生 Miss 且该组已满，则根据替换算法选择一个 Cache Line 换出。
\end{enumerate}

该过程体现了 Cache 的基本工作逻辑：先利用 Index 定位组，再在组内进行 Tag 匹配，最后根据匹配结果判断命中或缺失。

\subsubsection{替换算法实现}

当某个 Cache Set 已满且新的内存块需要调入时，模拟器需要根据替换算法选择被换出的 Cache Line。本实验主要实现了 FIFO 和 LRU 两种策略。

\paragraph{FIFO 替换算法}

FIFO 即先进先出策略。该算法记录每个 Cache Line 被装入 Cache 的时间，当需要替换时，选择最早进入 Cache 的块换出。FIFO 的实现简单，但它并不关心 Cache Line 最近是否被访问过，因此可能会把仍然频繁使用的数据块换出。

\paragraph{LRU 替换算法}

LRU 即最近最少使用策略。该算法记录每个 Cache Line 最近一次被访问的时间，当需要替换时，选择最久没有被访问的块换出。相比 FIFO，LRU 更符合时间局部性原理，通常能够更好地保留近期可能再次访问的数据。

\subsubsection{核心访问流程}

Cache 模拟器的核心访问逻辑可以概括如下：

\begin{lstlisting}[language=C, caption={Cache 访问模拟的核心逻辑}]
AccessCache(addr):
    block_addr = addr / block_size
    set_index  = block_addr % set_num
    tag        = block_addr / set_num

    for each line in cache[set_index]:
        if line.valid == true and line.tag == tag:
            hit++
            update_lru_information()
            return

    miss++

    if there is an empty line in cache[set_index]:
        fill the empty line with new tag
    else:
        victim = select_victim_by_policy()
        replace victim line with new tag
\end{lstlisting}

由此可见，Cache 模拟器的关键在于正确实现地址映射和替换逻辑。只要 trace.log 的访存地址能够被正确解析，模拟器就可以在不同 Cache 参数下重复运行，从而获得可比较的实验数据。